% =============================================================================
% 6. FUNDAMENTOS DE ARQUITECTURA
% =============================================================================

\section{Fundamentos de Arquitectura}

\subsection{Estilo Arquitectónico}

\subsubsection{Arquitectura de Marketplace Multinivel con Comunicación Orientada a Eventos}

El estilo arquitectónico adoptado por Smart-Park es una \textbf{arquitectura de marketplace multinivel con comunicación orientada a eventos}. Este estilo combina tres principios fundamentales:

\begin{enumerate}
    \item \textbf{Marketplace multi-tenant:} la plataforma aloja múltiples estacionamientos (inquilinos) bajo una misma infraestructura compartida, donde cada local conserva autonomía sobre sus datos, configuración, tarifas y operaciones. El aislamiento lógico se logra mediante identificadores de local en cada tabla y middleware de autorización que filtra los datos según el contexto del usuario autenticado (Richards y Ford, 2020).
    \item \textbf{Arquitectura en capas (\textit{layered architecture}):} el sistema se organiza verticalmente en cinco capas funcionales (presentación, aplicación, servicios locales, mensajería y caché, y datos), cada una con una responsabilidad claramente definida. La comunicación entre capas es descendente y unidireccional: cada capa solo depende de la capa inmediatamente inferior, lo que facilita la sustitución de componentes y la prueba unitaria (Bass et al., 2021).
    \item \textbf{Comunicación orientada a eventos (\textit{event-driven}):} las operaciones que no requieren respuesta inmediata (notificaciones, liberación de reservas, procesamiento de imágenes, copias de seguridad) se desacoplan mediante una cola de mensajes (RabbitMQ). Los productores publican eventos en intercambios (\textit{exchanges}) temáticos, y los consumidores los procesan de forma asíncrona, lo que mejora la escalabilidad y la tolerancia a fallos (Hohpe y Woolf, 2003).
\end{enumerate}

La justificación de este estilo se basa en los siguientes \textit{drivers} de arquitectura:

\begin{longtable}{p{3cm} p{11cm}}
\caption{Justificación del Estilo Arquitectónico} \\
\toprule
\textbf{Driver} & \textbf{Cómo lo satisface el estilo} \\
\midrule
\endfirsthead

\caption[]{Justificación del Estilo Arquitectónico (continuación)} \\
\toprule
\textbf{Driver} & \textbf{Cómo lo satisface el estilo} \\
\midrule
\endhead

\midrule
\multicolumn{2}{r}{\footnotesize\textit{Continúa en la siguiente página}} \\
\endfoot

\bottomrule
\endlastfoot

Escalabilidad & La separación en capas permite escalar horizontalmente la capa de aplicación (múltiples réplicas del contenedor) y verticalmente la capa de datos (PostgreSQL con mayor CPU/RAM). La cola de mensajes absorbe picos de carga sin afectar la API. \\
\midrule
Disponibilidad & Los procesos asíncronos desacoplados mediante RabbitMQ aseguran que un fallo en el motor de visión artificial o en el servicio de notificaciones no derrumbe la API central. \\
\midrule
Seguridad & La capa de aplicación centraliza la autenticación y autorización (RBAC + JWT), y el aislamiento multi-tenant impide el acceso cruzado entre locales. \\
\midrule
Mantenibilidad & La separación en capas y módulos permite que el equipo modifique un componente sin afectar los demás. Los contratos de interfase (JSON/REST) formalizan la comunicación. \\
\midrule
Rendimiento & Redis como caché reduce la latencia de las consultas frecuentes. WebSocket elimina la necesidad de \textit{polling} para la actualización en tiempo real. \\
\end{longtable}

\textit{Nota.} RBAC = Role-Based Access Control. JWT = JSON Web Token. Los drivers de arquitectura se detallan en la sección 8.5.2.

\subsubsection{Patrones Complementarios del Estilo Arquitectónico}

El estilo arquitectónico principal se complementa con los siguientes patrones:

\begin{longtable}{p{3.5cm} p{10.5cm}}
\caption{Patrones Complementarios del Estilo Arquitectónico} \\
\toprule
\textbf{Patrón} & \textbf{Aplicación en Smart-Park} \\
\midrule
\endfirsthead

\caption[]{Patrones Complementarios del Estilo Arquitectónico (continuación)} \\
\toprule
\textbf{Patrón} & \textbf{Aplicación en Smart-Park} \\
\midrule
\endhead

\midrule
\multicolumn{2}{r}{\footnotesize\textit{Continúa en la siguiente página}} \\
\endfoot

\bottomrule
\endlastfoot

API Gateway & La API central actúa como punto de entrada único para todos los clientes. Centraliza la autenticación, la autorización, la validación de datos y el enrutamiento a los módulos internos. \\
\midrule
Publish-Subscribe & Los eventos del sistema (reservas, notificaciones, detecciones de cámara) se publican en intercambios de RabbitMQ. Los consumidores suscritos procesan los eventos de forma independiente. \\
\midrule
CQRS (simplificado) & Las operaciones de escritura (crear reserva, registrar pago) se procesan de forma síncrona en la API, mientras que las operaciones de lectura frecuentes (disponibilidad, tarifas) se atienden desde el caché Redis. \\
\midrule
Cache-Aside & Antes de consultar PostgreSQL, la API verifica si el dato existe en Redis. Si existe, lo retorna directamente; si no, consulta la base de datos, almacena el resultado en Redis y lo retorna. \\
\midrule
Circuit Breaker & Las llamadas a servicios externos (pasarelas de pago, cámaras IP) implementan un \textit{circuit breaker} que interrumpe los reintentos cuando el servicio externo no responde, evitando la cascada de fallos. \\
\midrule
Observer (WebSocket) & El servidor de WebSocket implementa el patrón observador: los clientes se suscriben a canales por local, y el servidor notifica a todos los suscriptores cuando el estado de un espacio cambia. \\
\end{longtable}

\textit{Nota.} CQRS = Command Query Responsibility Segregation. Los patrones se basan en las referencias de Hohpe y Woolf (2003) y Gamma et al. (1994).

\subsection{Frameworks}

Los frameworks y bibliotecas utilizados en la construcción del sistema se seleccionaron considerando la madurez del ecosistema, la documentación disponible, la compatibilidad con las decisiones de arquitectura y la experiencia del equipo de desarrollo.

\begin{longtable}{p{3cm} p{2cm} p{9cm}}
\caption{Frameworks y Bibliotecas del Sistema} \\
\toprule
\textbf{Framework / Biblioteca} & \textbf{Capa} & \textbf{Justificación} \\
\midrule
\endfirsthead

\caption[]{Frameworks y Bibliotecas del Sistema (continuación)} \\
\toprule
\textbf{Framework / Biblioteca} & \textbf{Capa} & \textbf{Justificación} \\
\midrule
\endhead

\midrule
\multicolumn{3}{r}{\footnotesize\textit{Continúa en la siguiente página}} \\
\endfoot

\bottomrule
\endlastfoot

React 19 & Presentación & Biblioteca de componentes declarativos con DOM virtual, amplio ecosistema y soporte a largo plazo. Permite construir interfaces dinámicas e interactivas con estado reactivo (Meta Platforms, Inc., 2024). \\
\midrule
Vite 6 & Presentación & Empaquetador de módulos con recarga en caliente (HMR) instantánea, compilación rápida basada en esbuild y Rollup, y soporte nativo para TypeScript y JSX. \\
\midrule
React Router 7 & Presentación & Enrutador declarativo para SPA con rutas anidadas, rutas protegidas y carga diferida (\textit{lazy loading}). \\
\midrule
Leaflet + React-Leaflet & Presentación & Biblioteca de mapas interactivos de código abierto, ligera y extensible. Soporta capas de calles (OpenStreetMap) y satélite (Esri). \\
\midrule
FastAPI 0.115 & Aplicación & Framework web asíncrono de alto rendimiento para Python, con validación automática de datos (Pydantic), documentación OpenAPI generada y soporte nativo para WebSocket (Ramírez, 2018). \\
\midrule
SQLAlchemy 2.0 & Aplicación & ORM y toolkit SQL para Python con soporte asíncrono (\texttt{asyncio}), modelo declarativo y gestión de sesiones (SQLAlchemy, 2024). \\
\midrule
Pydantic 2.0 & Aplicación & Validación y serialización de datos mediante modelos tipados, con integración nativa en FastAPI. \\
\midrule
Uvicorn & Aplicación & Servidor ASGI de alto rendimiento basado en \texttt{uvloop} y \texttt{httptools}, compatible con FastAPI. \\
\midrule
YOLOv8 (Ultralytics) & Servicios locales & Modelo de detección de objetos en tiempo real, preentrenado en COCO, ajustable para la detección de vehículos (Ultralytics, 2024). \\
\midrule
OpenCV & Servicios locales & Biblioteca de visión artificial de código abierto, utilizada como respaldo y para el preprocesamiento de imágenes. \\
\midrule
RabbitMQ (Pika) & Mensajería & Cliente AMQP para Python que gestiona la conexión, los intercambios y las colas de mensajes (Pivotal Software, Inc., 2024). \\
\midrule
Redis (redis-py) & Caché & Cliente Redis asíncrono para Python, utilizado para operaciones de caché, sesiones y contadores (Redis Ltd., 2024). \\
\midrule
Docker & Infraestructura & Contenedorización del backend con imagen multi-stage para optimizar el tamaño y la seguridad del contenedor de producción (Docker, Inc., 2024). \\
\end{longtable}

\textit{Nota.} ORM = Object-Relational Mapping. ASGI = Asynchronous Server Gateway Interface. HMR = Hot Module Replacement. COCO = Common Objects in Context.

\subsection{Patrones de Diseño}

Los patrones de diseño aplicados en Smart-Park resuelven problemas recurrentes de diseño orientado a objetos y de arquitectura de software. Se clasifican según su propósito:

\begin{longtable}{p{2.5cm} p{2.5cm} p{9cm}}
\caption{Patrones de Diseño Aplicados} \\
\toprule
\textbf{Patrón} & \textbf{Categoría} & \textbf{Aplicación en Smart-Park} \\
\midrule
\endfirsthead

\caption[]{Patrones de Diseño Aplicados (continuación)} \\
\toprule
\textbf{Patrón} & \textbf{Categoría} & \textbf{Aplicación en Smart-Park} \\
\midrule
\endhead

\midrule
\multicolumn{3}{r}{\footnotesize\textit{Continúa en la siguiente página}} \\
\endfoot

\bottomrule
\endlastfoot

Repository & Estructural & Cada entidad del dominio (reservas, usuarios, estacionamientos, espacios) tiene un repositorio que encapsula las consultas a PostgreSQL mediante SQLAlchemy. La capa de servicios de negocio consume los repositorios sin conocer los detalles del ORM (Fowler, 2002). \\
\midrule
Service Layer & Estructural & La lógica de negocio se organiza en servicios (\texttt{ReservaService}, \texttt{PagoService}, \texttt{EspacioService}) que orquestan las operaciones entre repositorios, la cola de mensajes y los servicios externos. \\
\midrule
Factory Method & Creacional & La creación de reservas, pagos y notificaciones se delega a métodos de fábrica que encapsulan la lógica de instanciación según el tipo de modalidad, método de pago o canal de notificación. \\
\midrule
Strategy & Comportamiento & El cálculo de tarifas utiliza el patrón estrategia: cada tipo de tarifa (por hora, por minuto, abono, con recargo nocturno) implementa una interfaz común \texttt{TarifaStrategy}, y el servicio de reservas selecciona la estrategia adecuada según la configuración del local. \\
\midrule
Observer & Comportamiento & El servidor de WebSocket implementa el patrón observador. Los clientes se registran como observadores de un canal (identificado por el local), y el servidor les notifica cuando el estado de un espacio cambia. \\
\midrule
Singleton & Creacional & Las conexiones a la base de datos, a Redis y a RabbitMQ se gestionan como instancias únicas (\textit{singletons}) para evitar la proliferación de conexiones y garantizar la reutilización de recursos. \\
\midrule
Decorator & Estructural & Los middleware de FastAPI (\texttt{@require\_role}, \texttt{@rate\_limit}, \texttt{@audit\_log}) se implementan como decoradores que envuelven los endpoints con lógica transversal (autorización, limitación de intentos, auditoría). \\
\midrule
DTO & Estructural & Los modelos Pydantic (\texttt{ReservaCreate}, \texttt{ReservaResponse}, \texttt{PagoRequest}) actúan como Data Transfer Objects que validan y serializan los datos de entrada y salida de la API, separando la representación externa de las entidades internas del dominio. \\
\end{longtable}

\textit{Nota.} DTO = Data Transfer Object. ORM = Object-Relational Mapping. Los patrones se basan en las referencias de Gamma et al. (1994), Fowler (2002) y Martin (2017). La combinación de estos patrones garantiza una base de código mantenible, extensible y alineada con los principios SOLID.
